% ---- begin paper/tables/t1_descriptivos.tex
\begin{table}[htbp]\centering
\begin{minipage}{360pt}
\caption{Descriptive statistics.}
\label{tab:sample}
{\small
\begin{tabular*}{360pt}{@{}l@{\extracolsep{\fill}}lll@{}}
\toprule
 & (1) & (2) & (3) \\
 & Non-Earthquake & Earthquake & Difference \\
\midrule
Girl & 0.489 & 0.489 & $-$0.000 \\
 & (0.500) & (0.500) & (0.011) \\
Age at the earthquake (months) & 26.930 & 27.098 & 0.169 \\
 & (12.710) & (12.614) & (0.254) \\
Mother's age (2010) & 28.980 & 29.676 & 0.695\sym{***} \\
 & (7.066) & (7.116) & (0.196) \\
Mother's schooling (2010) & 10.312 & 10.311 & $-$0.002 \\
 & (3.182) & (3.180) & (0.211) \\
HH size (2010) & 4.865 & 4.948 & 0.083 \\
 & (1.641) & (1.686) & (0.074) \\
Rural household (2010) & 0.080 & 0.094 & 0.014 \\
 & (0.271) & (0.292) & (0.035) \\
Children of the mother (2012) & 2.165 & 2.155 & $-$0.010 \\
 & (1.072) & (1.075) & (0.029) \\
Prior mental health, mother (2012) & 0.091 & 0.125 & 0.034\sym{***} \\
 & (0.288) & (0.331) & (0.012) \\
Prior mental health, father (2012) & 0.023 & 0.027 & 0.004 \\
 & (0.150) & (0.161) & (0.004) \\
Prior mental health, relative (2012) & 0.123 & 0.152 & 0.029\sym{*} \\
 & (0.328) & (0.359) & (0.017) \\
\addlinespace[6pt]
Observations & 2,349 & 7,654 & 10,003 \\
\bottomrule
\end{tabular*}\par}
\vspace{4pt}
{\footnotesize
\textit{Notes}: ELPI uses sampling weight at national level. The sample are the 10,003 adolescents of the 2024 wave, born January 2006--August 2009, aged 6 to 49 months on the day of the earthquake and 15 to 18 at the 2024 interview. Sex and ages are fixed at birth and the 2010 rows come from the pre-earthquake baseline; the 2012 rows are the family size and the lifetime conditions reported in the 2012 wave, the controls of the main specification, and are missing for adolescents without that interview. Mental health rows mark a reported diagnosis of depression, anxiety, bipolar disorder, schizophrenia, a personality disorder or an autism spectrum disorder. Earthquake municipalities are those of the six regions where the official intensity study records destructive shaking (V, VI, VII, VIII, IX and Metropolitan; Astroza et al. 2010). Column 3 is the difference from a regression on the earthquake indicator; standard errors clustered at the municipality of selection appear in parentheses---number of clusters: 116. * p $<$ 0.1, ** p $<$ 0.05, *** p $<$ 0.01.
\par}
\end{minipage}
\end{table}
% ---- end paper/tables/t1_descriptivos.tex
